% =====================================================================
\part{Prüfung}
% =====================================================================

\chapter{Berichtete Ergebnisse}
\label{chap:berichtet}

Dieses Kapitel gibt wieder, was die README und der Abdeckungsbericht (\texttt{doc/coverage}, erzeugt am 21.~Oktober 2026, 11:54 Uhr, Pfade mit Rückstrich, also auf Windows) berichten, und verbindet es mit den Sätzen dieser Arbeit. Es wurde nichts neu ausgeführt.

\section{Umfang und Abdeckung}

Das Paket umfasst 14 Module mit $1\,170$ Zeilen, die Tests 11 Dateien mit $1\,150$ Zeilen; die README nennt 126 bestandene Tests und eine Abdeckung von 99\,\%. Im Quelltext der Tests stehen 85 Testfunktionen, davon 10 parametrisiert; die Differenz zu 126 ist mit der Parametrisierung verträglich, wurde aber nicht nachgerechnet. Tabelle~\ref{tab:abdeckung} gibt den Bericht wieder.

\begin{table}[!ht]
	\centering
	\small
	\begin{tabular}{@{}lrrrrr@{}}
		\toprule
		Modul & Anweisungen & fehlend & Zweige & teilweise & Abdeckung \\
		\midrule
		\texttt{\_\_init\_\_} & 9 & 0 & 0 & 0 & 100\,\% \\
		\texttt{\_\_main\_\_} & 3 & 0 & 0 & 0 & 100\,\% \\
		\texttt{adapters} & 65 & 0 & 0 & 0 & 100\,\% \\
		\texttt{binding} & 25 & 0 & 6 & 0 & 100\,\% \\
		\texttt{build} & 54 & 0 & 12 & 0 & 100\,\% \\
		\texttt{capi} & 41 & 0 & 2 & 0 & 100\,\% \\
		\texttt{cli} & 63 & 0 & 12 & 0 & 100\,\% \\
		\texttt{codegen} & 66 & 0 & 20 & 0 & 100\,\% \\
		\texttt{container} & 77 & 0 & 22 & 0 & 100\,\% \\
		\texttt{errors} & 27 & 0 & 0 & 0 & 100\,\% \\
		\texttt{loader} & 135 & 0 & 30 & 0 & 100\,\% \\
		\texttt{metadata} & 99 & 0 & 40 & 0 & 100\,\% \\
		\texttt{platform\_info} & 21 & 0 & 6 & 0 & 100\,\% \\
		\texttt{search} & 95 & 4 & 50 & 4 & 96\,\% \\
		\midrule
		Summe & 780 & 4 & 200 & 4 & 99\,\% \\
		\bottomrule
	\end{tabular}
	\caption{Abdeckungsbericht (Berichtet). Zweigabdeckung ist eingeschaltet (\texttt{branch = true} in \texttt{pyproject.toml}).}
	\label{tab:abdeckung}
\end{table}

Die vier fehlenden Anweisungen und vier teilweise abgedeckten Zweige liegen alle in \texttt{search.py}: Zeile~21 (Rückgabe im Semikolon-Zweig), die Zeilen~125 bis 127 (POSIX-Zweig von \texttt{\_world\_writable}, auf Windows nicht erreicht), sowie die Zweige 20$\to$21, 60$\to$65, 62$\to$65 und 120$\to$125.

\begin{lemma}[Ein toter Zweig in \texttt{split\_sl\_path}]
	\label{lem:tot}
	Im Listing~\ref{lst:split} ist der Zweig „\texttt{if token} ist falsch“ am Ende der Funktion (Sprung 62$\to$65) unerreichbar. Der Zweig „\texttt{if current} ist falsch“ (Sprung 60$\to$65) ist erreichbar, etwa für \texttt{"/a:"}.
\end{lemma}

\begin{proof}
	\texttt{current} wird nur durch \texttt{append} mit einem Ein-Zeichen-String (\texttt{raw[i]}, \texttt{":"} oder \texttt{"\textbackslash\textbackslash"}) gefüllt; ist die Liste nichtleer, ist \texttt{"".join(current)} ein nichtleerer String, also ist \texttt{token} wahr. Der Zweig 62$\to$65 setzt dagegen ein falsches \texttt{token} bei wahrem \texttt{current} voraus. Für \texttt{"/a:"} wird beim letzten Doppelpunkt \texttt{current} auf die leere Liste gesetzt, die Schleife endet, und \texttt{if current} ist falsch.
\end{proof}

Der tote Zweig ist harmlos, aber auch nicht abdeckbar; ein Test kann ihn nicht erreichen. Die Prüfung \texttt{if token} im Schlussteil kann entfallen; das vereinfacht die Funktion und hebt die Zweigabdeckung der Datei.

\section{Welche Sätze die Tests belegen}

Tabelle~\ref{tab:testbeleg} ordnet den Sätzen dieser Arbeit die Tests zu, die ihre Aussage an Beispielen prüfen, und nennt die Lücken. Ein Test ersetzt keinen Beweis, bindet ihn aber an den Quelltext: Ändert jemand das Verhalten, schlägt der Test an.

{\small
\begin{longtable}{@{}L{4.3cm}L{6.5cm}L{3.4cm}@{}}
\caption{Beleg der Sätze durch Tests (Namen aus \texttt{tests/}).}\label{tab:testbeleg}\\
\toprule
Aussage & Test & Lücke \\
\midrule
\endfirsthead
\toprule
Aussage & Test & Lücke \\
\midrule
\endhead
\bottomrule
\endfoot
Satz~\ref{satz:rundreise} Rundreise & \texttt{test\_fat\_roundtrip}, \texttt{test\_roundtrip\_and\_with\_targets} & keine Zufallseingaben \\[0.3em]
Lemma~\ref{lem:akzept} Prüfungen & \texttt{test\_short\_magic\_version\_and\_truncation}, \texttt{test\_unreadable\_metadata}, \texttt{test\_table\_errors} & \\[0.3em]
Satz~\ref{satz:integritaet} Integrität & \texttt{test\_hash\_mismatch} & Metadaten ungeschützt (nicht geprüft) \\[0.3em]
Satz~\ref{satz:kanon} Nichtkanonizität & -- & kein Test \\[0.3em]
Lemma~\ref{lem:name} Namen & \texttt{test\_resolve\_name} & \\[0.3em]
Satz~\ref{satz:struktur} Struktur & \texttt{test\_candidate\_order\_and\_filters} & Windows-Systemverzeichnisse \\[0.3em]
Satz~\ref{satz:cwd} Arbeitsverzeichnis & \texttt{test\_cwd\_is\_never\_searched} & relatives $R$ \\[0.3em]
Satz~\ref{satz:auswahl} erster Treffer & \texttt{test\_locate\_hits\_by\_priority} & \\[0.3em]
Satz~\ref{satz:gehaertet} gehärtet & \texttt{test\_hardened\_rejects\_world\_writable}, \texttt{test\_requester\_dir\_and\_hardened} & Gruppenrechte, Elternverzeichnisse \\[0.3em]
Satz~\ref{satz:split} \texttt{SL\_PATH} & \texttt{test\_split\_sl\_path} & Semikolon-Zweig, Laufwerke \\[0.3em]
Satz~\ref{satz:erfolg} Zyklen & \texttt{test\_cycle\_is\_detected} & \\[0.3em]
Lemma~\ref{lem:vertrag}(a) Rollback & \texttt{test\_missing\_dependency\_rolls\_back}, \texttt{test\_init\_failure\_rolls\_back}, \texttt{test\_mapping\_failure} & Nicht-\texttt{SLError} \\[0.3em]
Satz~\ref{satz:refcount} Zählung & \texttt{test\_refcount\_shares\_handle}, \texttt{test\_shared\_dependency\_is\_refcounted}, \texttt{test\_closed\_handle\_is\_rejected} & \\[0.3em]
Satz~\ref{satz:reihenfolge} Reihenfolge & \texttt{test\_dependency\_chain\_and\_release}, \texttt{test\_init\_runs\_once\_and\_fini\_on\_close}, \texttt{test\_missing\_fini\_does\_not\_block\_unload} & \\[0.3em]
Satz~\ref{satz:export} Export & \texttt{test\_only\_exported\_symbols\_are\_visible}, \texttt{test\_real\_library\_roundtrip} & \\[0.3em]
Satz~\ref{satz:deadlock} Deadlock & -- & kein Test \\[0.3em]
Satz~\ref{satz:schichtkonsistenz} Schichten & \texttt{test\_generate\_all}, \texttt{test\_python\_binding\_is\_valid\_and\_loadable}, \texttt{test\_generated\_python\_binding\_against\_real\_library} & \\[0.3em]
Lemma~\ref{lem:parse}, Satz~\ref{satz:typen} & \texttt{test\_parse\_signature}, \texttt{test\_parse\_signature\_rejects}, \texttt{test\_ctype\_expr} & \\[0.3em]
Satz~\ref{satz:build} Build & \texttt{test\_build\_sl\_adds\_host\_to\_extra\_payloads}, \texttt{test\_fat\_container\_and\_cli\_build} & \\[0.3em]
Satz~\ref{satz:konvex} Bereiche & \texttt{test\_satisfies}, \texttt{test\_satisfies\_rejects\_bad\_range} & unerfüllbare Bereiche \\
\end{longtable}
}

Drei Lücken fallen auf: Es gibt keinen Test für den Semikolon-Zweig von \texttt{SL\_PATH} (der Bericht weist ihn als nie ausgeführt aus), keinen für Ausnahmen außerhalb von \texttt{SLError} in Hooks und keinen für nebenläufige Aufrufe. Alle drei sind in den Befunden \ref{bef:sep}, \ref{bef:exc} und \ref{bef:lock} wieder aufgenommen.

\section{Einordnung der README-Aussagen}

Die README sagt, geprüft worden sei unter Linux, und die Zeilenabdeckung habe bei 100\,\% gelegen. Der Bericht ist auf Windows erzeugt und weist für \texttt{search.py} 96\,\% aus (Tabelle~\ref{tab:abdeckung}). Die beiden Aussagen beschreiben vermutlich verschiedene Läufe; welcher gilt, lässt sich ohne Neulauf nicht sagen \offen{nicht geprüft}. Die Unterschiede sind in Befund~\ref{bef:doku} festgehalten.

\chapter{Befunde und Verbesserungsvorschläge}
\label{chap:befunde}

Dieses Kapitel fasst vierzehn Befunde zusammen. Jeder hat eine Kennung \textbf{B1} bis \textbf{B14}, einen Beleg (Satz oder Quelltextstelle) und einen Vorschlag. Befunde sind keine Fehler im strengen Sinn: Teils weichen sie vom Entwurf ab, teils benennen sie Annahmen, unter denen die Beweise gelten.

\begin{befund}[B1: \texttt{SL\_PATH} mit zwei Betriebsarten]
	\label{bef:sep}
	\emph{Beleg:} Satz~\ref{satz:split}. Die Zerlegung wählt nach dem Inhalt (Semikolon vorhanden) die Maskierungsregeln; Einträge mit Semikolon sind nicht darstellbar; der Laufwerkszweig wirkt nur für den ersten Eintrag; der Semikolon-Zweig ist im Bericht ungetestet. \emph{Vorschlag:} ein einziges Trennzeichen und eine Maskierung ohne Sonderfälle (Abschnitt~\ref{sec:slpath-ausblick}), dazu Tests für beide Fälle und die Dokumentation in README und Entwurf.
\end{befund}

\begin{befund}[B2: Relative Verzeichnisse]
	\label{bef:cwd}
	\emph{Beleg:} Satz~\ref{satz:cwd}. \texttt{requester\_dir} und \texttt{app\_dir} werden nicht normalisiert; unter Windows können \texttt{ProgramData} und \texttt{ProgramFiles} relativ sein. Dann hängt die Suche vom Arbeitsverzeichnis ab. \emph{Vorschlag:}
\begin{lstlisting}[language=Python]
self.requester_dir = os.path.abspath(requester_dir) if requester_dir else None
self.app_dir = os.path.abspath(app_dir) if app_dir else None
# und im WindowsAdapter: nur absolute Systemverzeichnisse verwenden
\end{lstlisting}
\end{befund}

\begin{befund}[B3: Rechteprüfung im gehärteten Modus]
	\label{bef:ww}
	\emph{Beleg:} Satz~\ref{satz:gehaertet} und \texttt{search.py} Zeilen 113--127. Geprüft wird nur das Bit \emph{others} des unmittelbaren Verzeichnisses; bei \texttt{OSError} wird zugelassen; unter Windows entscheidet der Verzeichnisname \texttt{open} (ein Testname) über das Ergebnis. \emph{Vorschlag:} Besitzer und Rechte aller übergeordneten Verzeichnisse prüfen, bei Fehlern ablehnen, Windows-ACLs über eine eigene Funktion des Adapters abfragen, und den Testnamen aus dem Produktivcode entfernen (der Test kann sein Verzeichnis über einen Adapter-Stub kennzeichnen).
\end{befund}

\begin{befund}[B4: Globale Sperre über Hooks]
	\label{bef:lock}
	\emph{Beleg:} Satz~\ref{satz:deadlock}. Der Entwurf verlangt Initialisierung ohne eine Sperre, die Nutzercode betreten kann. \emph{Vorschlag:} Sperren je Schlüssel mit Zuständen \emph{Laden}/\emph{Geladen} und Aufruf der Hooks außerhalb jeder Lader-Sperre (Abschnitt~\ref{sec:schluessel-sperren}), dazu ein Test mit zwei Threads.
\end{befund}

\begin{befund}[B5: Ausnahmen außerhalb von \texttt{SLError}]
	\label{bef:exc}
	\emph{Beleg:} Satz~\ref{satz:ausnahme}(b). Unter Windows löst ein Speicherzugriffsfehler in nativem Code in \texttt{ctypes} eine \texttt{OSError} aus, die den Rollback umgeht. \emph{Vorschlag:}
\begin{lstlisting}[language=Python]
    except BaseException:
        if native is not None:
            self.adapter.unmap(native)
        for dep_handle in deps:
            self._release(dep_handle)
        raise
\end{lstlisting}
	und im Fini-Pfad \texttt{unmap} und die Rückgabe der Abhängigkeiten in ein \texttt{finally}.
\end{befund}

\begin{befund}[B6: Native Abhängigkeiten]
	\label{bef:native}
	\emph{Beleg:} Quelltext \texttt{adapters.py} Zeile 40, \texttt{ctypes.CDLL(path)} ohne Modus. Nach der Dokumentation von \texttt{ctypes} ist der Vorgabemodus unter POSIX \emph{lokal}. Dann stellen die Metadaten-Abhängigkeiten die Reihenfolge von Init und Freigabe sicher, machen die Symbole der Abhängigkeit aber nicht für die abhängige Bibliothek auflösbar. Der Test \texttt{test\_real\_dependency\_and\_errors} deklariert eine Abhängigkeit, die der Code nicht nutzt \offen{Wirkung nicht ausgeführt geprüft}. \emph{Vorschlag:} Abbildungsmodus je Abhängigkeit in den Metadaten angeben (global/lokal) oder einen eigenen Abbilder einsetzen (Abschnitt~\ref{sec:runtime-ausblick}), und einen Test, in dem die abhängige Bibliothek ein Symbol der Basis aufruft.
\end{befund}

\begin{befund}[B7: Container nicht kanonisch, ohne Echtheit]
	\label{bef:canon}
	\emph{Beleg:} Satz~\ref{satz:kanon}, Bemerkung~\ref{bem:hashgrenze}. \emph{Vorschlag:} kanonische Kodierung erzwingen (keine überzähligen Bytes, geordnete Tabelle, disjunkte Bereiche, kanonisches JSON) und eine Signatur über einen Digest von Metadaten, Tabelle und Payload-Hashes (Abschnitt~\ref{sec:signatur-ausblick}).
\end{befund}

\begin{befund}[B8: Eingabeprüfung der Metadaten]
	\label{bef:parse}
	\emph{Beleg:} \texttt{metadata.py} Zeile 137: \texttt{tuple(data.get("exports", ()))} macht aus einer Zeichenkette \texttt{"wsat"} das Tupel \texttt{('w','s','a','t')}, das als Liste gültiger Ein-Zeichen-Bezeichner akzeptiert wird. Außerdem fängt \texttt{unpack} bei der Metadaten-Zerlegung nur \texttt{ValueError}; sehr tief verschachtelte JSON-Eingaben lösen in CPython eine \texttt{RecursionError} aus \offen{nicht ausgeführt}. \emph{Vorschlag:} Typprüfung (\texttt{isinstance(value, list)}) für Listenfelder und \texttt{except (ValueError, RecursionError)}.
\end{befund}

\begin{befund}[B9: Build-Werkzeug]
	\label{bef:build}
	\emph{Beleg:} Satz~\ref{satz:build} und die Folgetexte. GNU-Flags für alle Windows-Compiler, auch für \texttt{cl}; keine Angabe von Sprachstandard, Bibliotheken oder Linker; \texttt{cc} linkt keine C++-Laufzeit. \emph{Vorschlag:} Optionen \texttt{--cc}, \texttt{--std}, \texttt{-l}, \texttt{--ldflag} und getrennte Toolchain-Beschreibungen für GNU und MSVC (Aufruf \texttt{cl /LD /Fe:}).
\end{befund}

\begin{befund}[B10: Erzeugtes CMake-Modul]
	\label{bef:cmake}
	\emph{Beleg:} \texttt{codegen.py} Zeilen 96--110. Das Modul benutzt \texttt{SL\_TARGET\_PLATFORM}, \texttt{SL\_TARGET\_ARCH} und \texttt{Python3\_EXECUTABLE}, ohne sie zu setzen, und stellt nur \texttt{C\_VISIBILITY\_PRESET} ein, obwohl die Zielprojekte C++ verwenden. \emph{Vorschlag:} die Variablen aus \texttt{CMAKE\_SYSTEM\_NAME} und \texttt{CMAKE\_SYSTEM\_PROCESSOR} ableiten, \texttt{find\_package(Python3)} aufrufen und \texttt{CXX\_VISIBILITY\_PRESET} und \texttt{VISIBILITY\_INLINES\_HIDDEN} ergänzen.
\end{befund}

\begin{befund}[B11: Felder und Daten ohne Wirkung]
	\label{bef:dead}
	\emph{Beleg:} Das Feld \texttt{runtime} wird gespeichert und validiert, aber vom Lader nirgends ausgewertet; der Parameter \texttt{flags} von \texttt{open} wird nicht benutzt; die erzeugte C-Quelle legt eine Tabelle \texttt{<name>\_sl\_exports} an, die nirgends verwendet wird (bei \texttt{-Wall} kann das eine Warnung auslösen \offen{nicht geprüft}); Lemma~\ref{lem:tot} nennt einen unerreichbaren Zweig. \emph{Vorschlag:} die Felder mit Semantik belegen oder streichen; die Tabelle als Selbsttest verwenden (Vergleich mit den Metadaten) oder entfernen.
\end{befund}

\begin{befund}[B12: Dokumentationsstand]
	\label{bef:doku}
	\emph{Beleg:} Kapitel~\ref{chap:berichtet}. README: „Geprüft wurde unter Linux“, „Zeilenabdeckung 100\,\%“, Trennzeichen „\texttt{:} unter POSIX oder \texttt{;} unter Windows“; Bericht: Windows-Lauf, 96\,\% in \texttt{search.py}; Code: Betriebsart nach Inhalt. \emph{Vorschlag:} Lauf und Bericht in CI erzeugen und die README-Angaben daraus ableiten.
\end{befund}

\begin{befund}[B13: Geteilte Handles]
	\label{bef:handle}
	\emph{Beleg:} Satz~\ref{satz:nutzung} und die Bemerkung danach. Aufrufer teilen ein Handle-Objekt; ein doppeltes \texttt{close} nimmt die Referenz eines anderen zurück; \texttt{symbol} und \texttt{info} nehmen die Sperre nicht. \emph{Vorschlag:} je \texttt{open} ein eigenes Zeugenobjekt (Token) herausgeben, das beim ersten \texttt{close} verbraucht wird, und \texttt{symbol} unter der Sperre des Schlüssels ausführen.
\end{befund}

\begin{befund}[B14: Kleinere Abweichungen vom Entwurf]
	\label{bef:abw}
	\emph{Beleg:} (a) \texttt{lib}-Präfix im logischen Namen zulässig (Lemma~\ref{lem:name}); (b) die Liste der geprüften Kandidaten wird bei Erfolg verworfen (\texttt{\_tried}); (c) \texttt{SLError.\_\_str\_\_} gibt nur Code und Meldung aus, nicht den mehrzeiligen Text des Entwurfs mit \texttt{failed\_name} und Suchliste; (d) Versionsbereiche brauchen dreiteilige Versionen (\texttt{>=1.2} ist ungültig); (e) das Anwendungsverzeichnis hat keinen Vorgabewert; (f) Konstruktoren und \texttt{DllMain} laufen im Systemlader beim Abbilden, nicht umhüllt; (g) kein Legacy-Modus für \texttt{.dll}/\texttt{.so}/\texttt{.dylib}; (h) \texttt{sl pack} und \texttt{sl build} ersetzen das Targetfeld der Metadaten (Satz~\ref{satz:build}); (i) die Lader-API ist Python, nicht die C-Schnittstelle des Entwurfs; (j) die Abbildung über eine temporäre Datei setzt ein Verzeichnis mit Ausführungsrecht voraus. \emph{Vorschlag:} die Punkte (a) bis (j) in die Spezifikation aufnehmen (als Abweichung begründen) oder umsetzen.
\end{befund}

Die Befunde B1, B2, B5 und B8 sind Änderungen von wenigen Zeilen; B4, B6 und B7 verlangen einen Entwurf, der im Ausblick skizziert wird.

% =====================================================================
\part{Schluss}
% =====================================================================

\chapter{Zusammenfassung}
\label{chap:zusammenfassung}

Die Arbeit hat das Kapitel zum Laden der C/C++-Bibliothek aus der Arbeit zur Wiederverwendung des Subgraph Algorithmus zu einer eigenständigen, formalen Untersuchung des Vorschlags eines einheitlichen Shared-Library-Formats \SL\ ausgebaut. Gegenstand sind der Diskussionsentwurf und die Referenzimplementierung \texttt{sharedlib}, die gelesen, aber nicht ausgeführt wurde. Die Antworten auf die acht Fragen der Einleitung lauten:

\begin{enumerate}
	\item \textbf{Plattformunterschiede.} Windows, Linux und macOS unterscheiden sich in Endung, Präfix, Binärformat, Lade-API, Standardsichtbarkeit, Abhängigkeitssuche, Initialisierung, Fehlermeldung und Entladen (Tabelle~\ref{tab:plattformen}). In den zwei Beispielprojekten sind sechs Stellen reiner Plattformcode. Dass der Maschinencode selbst nicht vereinheitlicht werden kann, ist keine Vorsicht des Entwurfs, sondern ein Satz: Kein Register-Code-Block bedient zwei Aufrufkonventionen (Satz~\ref{satz:nichtziel}).
	\item \textbf{Betriebssystemunabhängigkeit.} Sie lässt sich als Adapter-Parametrizität formulieren (Definition~\ref{def:parametrizitaet}). Die Referenzimplementierung hat sie bis auf fünf benannte Ausnahmen (Satz~\ref{satz:parametrisch}, Tabelle~\ref{tab:leaks}): Pfadfunktionen, die Rechteprüfung, \texttt{realpath} und die Vorgabewerte. Die Konformitätsmatrix (Tabelle~\ref{tab:konformitaet}) bewertet 25 Forderungen des Entwurfs; die meisten sind erfüllt, wenige teilweise, eine nicht, und eine fehlt.
	\item \textbf{Container.} Die Kodierung hat eine exakte Größenformel $18+m+80k+\sum|P(t)|$ (Lemma~\ref{lem:groesse}), eine bewiesene Rundreise (Satz~\ref{satz:rundreise}) und prüft sieben Bedingungen (Lemma~\ref{lem:akzept}). Der Hash schützt Payloads nur gegen Zufall (Satz~\ref{satz:integritaet}, Bemerkung~\ref{bem:hashgrenze}); der Container ist nicht kanonisch (Satz~\ref{satz:kanon}).
	\item \textbf{Suchordnung.} Die Kandidatenliste hat eine feste Struktur und hängt von der Umgebung nur über \texttt{SL\_PATH} und, auf Windows, zwei Systemvariablen ab (Satz~\ref{satz:struktur}); sie ist unabhängig vom Arbeitsverzeichnis, wenn alle konfigurierten Verzeichnisse absolut sind (Satz~\ref{satz:cwd}); der erste Treffer gewinnt (Satz~\ref{satz:auswahl}); der gehärtete Modus schließt \texttt{SL\_PATH} aus und lehnt beschreibbare Verzeichnisse ab (Satz~\ref{satz:gehaertet}). Die Zerlegung von \texttt{SL\_PATH} hat zwei vom Inhalt gewählte Betriebsarten (Satz~\ref{satz:split}).
	\item \textbf{Lader.} Erfolg ist durch Zyklenfreiheit sowie lokale und Kantenbedingungen charakterisiert (Satz~\ref{satz:erfolg}); Init und Freigabe laufen in Abhängigkeits- und umgekehrter Reihenfolge (Satz~\ref{satz:reihenfolge}); die Referenzzählung hat eine bewiesene Invariante (Satz~\ref{satz:refcount}); der Rollback stellt den Zustand wieder her, solange alle Ausnahmen \texttt{SLError} sind (Satz~\ref{satz:ausnahme}); die globale Sperre erlaubt einen Deadlock (Satz~\ref{satz:deadlock}).
	\item \textbf{Schichtung.} Der Importgraph ist azyklisch in sieben Ebenen, und der Laufzeitpfad schließt Codegenerierung, Build und Kommandozeile aus (Satz~\ref{satz:schichten}).
	\item \textbf{Codegenerierung und Build.} Alle Schichten sind reine Funktionen der Metadaten (Satz~\ref{satz:schichtkonsistenz}); die Typabbildung ist unter einer Annahme über \texttt{ctypes} ABI-treu (Satz~\ref{satz:typen}); das Build-Werkzeug ersetzt das Targetfeld (Satz~\ref{satz:build}) und deckt MSVC und C++ nicht ab.
	\item \textbf{Versionen.} Versionsbereiche sind Intervalle (Satz~\ref{satz:konvex}); die Auswahl unter mehreren Kandidaten mit Rückverfolgung ist NP-vollständig (Satz~\ref{satz:npvoll}); der Lader bleibt in Polynomialzeit, weil er den ersten Treffer nimmt, und ist deshalb unvollständig (Satz~\ref{satz:unvollstaendig}).
\end{enumerate}

Hinzu kommen eine Sicherheitsanalyse mit dem Lemma vom einmaligen Lesen (Lemma~\ref{lem:einmal}), die Auswertung des Abdeckungsberichts mit einem bewiesenen toten Zweig (Lemma~\ref{lem:tot}) und vierzehn Befunde. Ihr Kern: Der Entwurf trägt, die Referenzimplementierung setzt seine Zusagen für Container, Suche, Referenzzählung und Reihenfolgen nachweislich um; offen sind vor allem die Nebenläufigkeit (B4), die Ausnahmesicherheit (B5), native Abhängigkeiten (B6) und die Echtheit des Containers (B7). Für die Beispielprojekte folgt daraus ein ernüchterndes, aber klares Bild: Von den sechs Plattformstellen entfallen zwei ganz, drei teilweise, eine gar nicht (Tabelle~\ref{tab:ersatz}).

\chapter{Ausblick}
\label{chap:ausblick}

Der Ausblick ist nach Gegenständen gegliedert: zuerst die Behebung der Befunde, dann die Bausteine, die der Entwurf verlangt und die Referenzimplementierung noch nicht hat, schließlich die Verbindung zu den übrigen Projekten, Messungen, Standardisierung und offene Forschungsfragen. Wo möglich, ist die Aussage so formuliert, dass sie sich beweisen oder messen lässt; bewiesen ist davon \emph{nichts}, es sind Entwürfe und Beweisskizzen, und sie sind so gekennzeichnet.

\section{Die Befunde beheben: eine Reihenfolge}

Tabelle~\ref{tab:fahrplan} ordnet die Befunde nach Aufwand und Wirkung. Die Reihenfolge ist so gewählt, dass jede Stufe die Aussagen der Sätze nicht verschlechtert und jede Behebung einen Test erhält, der die Lücke aus Tabelle~\ref{tab:testbeleg} schließt.

\begin{table}[!ht]
	\centering
	\small
	\begin{tabularx}{\textwidth}{@{}L{1.2cm}L{3.0cm}YL{3.2cm}@{}}
		\toprule
		Stufe & Befunde & Inhalt & Test \\
		\midrule
		1 & B2, B5, B8, B11 & wenige Zeilen: Normalisierung, \texttt{BaseException}, Typprüfung, toter Code & Windows-Fall, Hook mit \texttt{OSError}, Container mit \texttt{exports} als Text \\[0.3em]
		2 & B1, B12, B14(a--e) & \texttt{SL\_PATH}-Grammatik, Dokumentation aus CI, Diagnosetext & Round-Trip über beide Betriebsarten \\[0.3em]
		3 & B3, B9, B10 & Rechteprüfung, Build-Optionen, CMake-Modul & Verzeichnisketten, Compiler-Matrix \\[0.3em]
		4 & B4, B13 & Sperren je Schlüssel, Token-Handles & Zwei-Thread-Test mit blockierendem Init \\[0.3em]
		5 & B6, B7 & native Abhängigkeiten, kanonischer Container mit Signatur & Aufruf in die Basis, Signaturvektoren \\
		\bottomrule
	\end{tabularx}
	\caption{Vorschlag für die Reihenfolge der Behebung.}
	\label{tab:fahrplan}
\end{table}

\section{Aussagen als Eigenschaftstests und mechanisch geprüfte Beweise}
\label{sec:verifikation-ausblick}

\subsection{Eigenschaftstests}

Die Tests der Referenzimplementierung sind Beispieltests. Fast jeder Satz der Arbeit lässt sich als Eigenschaft über zufällig erzeugte Eingaben formulieren und mit einem Werkzeug im Stil von QuickCheck \cite{claessen2000} prüfen:

\begin{itemize}
	\item \emph{Rundreise} (Satz~\ref{satz:rundreise}): für zufällige gültige $(M,P)$ gilt \texttt{unpack(pack($M,P$))} $=(M,P)$.
	\item \emph{Maskierung} (Satz~\ref{satz:split}(c)): für zufällige Listen nichtleerer, semikolonfreier Zeichenketten gilt $\mathrm{split}(\mathrm{join}(\cdot))=\mathrm{id}$ unter der Randbedingung.
	\item \emph{Zyklen und Erfolg} (Satz~\ref{satz:erfolg}): für zufällige Abhängigkeitsgraphen auf $\le 8$ Knoten (aus Fake-Containern gebaut) gelingt \texttt{open} genau dann, wenn der erreichbare Teilgraph azyklisch und lokal gültig ist; das Orakel dazu ist ein Algorithmus nach Tarjan \cite{tarjan1972}.
	\item \emph{Bilanz} (Satz~\ref{satz:refcount}): für zufällige Folgen von \texttt{open}/\texttt{close} ist am Ende $\Lambda=\emptyset$ und jedes \texttt{map} durch ein \texttt{unmap} beantwortet.
	\item \emph{Reihenfolge} (Satz~\ref{satz:reihenfolge}): in jedem Ablauf eines zufälligen Graphen ist die Folge der Init-Ereignisse eine umgekehrte topologische Ordnung.
	\item \emph{Rollback} (Lemma~\ref{lem:vertrag}): für zufällig fehlschlagende Orakel ist nach einem Fehler der Zustand gleich dem vorherigen.
	\item \emph{Struktur der Suche} (Satz~\ref{satz:struktur}): die Kandidatenliste ändert sich nicht, wenn \texttt{PATH}, \texttt{LD\_LIBRARY\_PATH} und \texttt{DYLD\_LIBRARY\_PATH} beliebig gesetzt werden.
\end{itemize}

Diese Tests sind klein, und sie fangen genau die Fehlerklasse, die Beispieltests übersehen: Randfälle der Verbindung von Eigenschaften. Der Test zum Beispiel~\ref{bsp:zweilaufwerke} würde B1 sofort zeigen.

\subsection{Mechanisch geprüfte Beweise}

Das Modell aus Kapitel~\ref{chap:grundlagen} ist klein genug für einen Beweisassistenten. Ein Vorgehen in drei Stufen:
\begin{enumerate}
	\item \emph{Zustandsmodell in TLA+} \cite{lamport2002}: Zustand $(\Lambda,\sigma)$, Aktionen \texttt{Open}, \texttt{Close}, Hooks, Orakelantworten als nichtdeterministische Wahl; Invarianten sind Satz~\ref{satz:refcount} und Abgeschlossenheit; Eigenschaften der Lebendigkeit sind der Ausschluss von Deadlocks in der Variante mit Schlüssel-Sperren. Der Modellprüfer wird auf Graphen bis zu fünf Knoten mit zwei Threads angesetzt.
	\item \emph{Beweise in einem Beweisassistenten} (Lean oder Coq): Satz~\ref{satz:erfolg} und Lemma~\ref{lem:vertrag} als Programmeigenschaften eines funktionalen Modells von \texttt{\_open}; Satz~\ref{satz:npvoll} als Reduktion mit Korrektheitsbeweis.
	\item \emph{Verbindung zum Quelltext:} Aus dem Python-Code wird ein funktionales Modell von Hand abgeleitet; Abweichungen zwischen Modell und Code sind die eigentliche Fehlerquelle, die Eigenschaftstests aus dem vorigen Abschnitt fangen.
\end{enumerate}
Die Kosten sind erheblich; der Nutzen besteht vor allem darin, dass die Annahmen (\ref{ann:determinismus}, \ref{ann:exc}, \ref{ann:ctypes}) sichtbar werden und jeder Satz nur diejenigen verwendet, die er braucht.

\section{Ein Lader mit Sperren je Schlüssel}
\label{sec:schluessel-sperren}

Befund~B4 verlangt, die globale Sperre durch eine feinere Steuerung zu ersetzen. Der folgende Entwurf ist nicht implementiert.

\subsection{Zustände und Regeln}

Jeder Schlüssel $k$ hat einen Zustand $z(k)\in\{\mathrm{Leer},\text{Lädt}(T),\mathrm{Geladen}\}$ mit dem ladenden Thread $T$. Eine kurze Sperre $\mathcal{L}$ schützt nur die Tabelle $z$ und die Zähler und wird \emph{nie} über einen Hook, ein \texttt{dlopen} oder einen Dateizugriff gehalten. Die Regeln:
\begin{enumerate}
	\item \texttt{open}$(k)$ setzt unter $\mathcal{L}$ den Zustand von \emph{Leer} auf \emph{Lädt}$(T)$ und lädt \emph{außerhalb} von $\mathcal{L}$; bei \emph{Geladen} wird der Zähler erhöht; bei \emph{Lädt}$(T')$ mit $T'\ne T$ wartet der Thread auf eine Bedingungsvariable zu $k$.
	\item Bei \emph{Lädt}$(T)$ durch denselben Thread (Wiedereintritt aus einem Hook) wird \texttt{CYCLE} gemeldet.
	\item Wartet $T$ auf $k$, wird vorher geprüft, ob das Warten einen Kreis im Warte-Graphen der Threads schlösse (Thread $T$ wartet auf $T'$, der auf $T''$ wartet \ldots auf $T$); in dem Fall wird \texttt{CYCLE} gemeldet.
	\item Hooks laufen ohne jede Lader-Sperre.
\end{enumerate}

\begin{satz}[Beweisskizze: Verklemmungsfreiheit des Entwurfs]
	\label{satz:skizze-deadlock}
	Im Entwurf kann kein Deadlock auftreten, der durch den Lader entsteht: Jeder Kreis im Warte-Graphen wird entweder durch Regel~3 verhindert, oder er enthält ein Warten, das nicht vom Lader ausgeht (ein Hook wartet auf Nutzercode).
\end{satz}

\begin{proof}[Skizze]
	Ein Thread wartet nur an Regel~1 (auf eine Bedingungsvariable zu $k$, die der ladende Thread $T'$ nach Abschluss oder Fehler signalisiert) und niemals auf $\mathcal{L}$ länger als für eine kurze Tabellenoperation. Ein Warten von $T$ auf $T'$ entsteht nur, wenn $T'$ einen Schlüssel lädt, den $T$ für eine Abhängigkeit braucht. Ein Kreis $T_1\to T_2\to\cdots\to T_1$ im Warte-Graphen bedeutete, dass jeder $T_i$ einen Schlüssel lädt, den $T_{i-1}$ braucht; die Wartekanten wären also Abhängigkeitskanten, und Regel~3 erkennt den Kreis beim Eintragen der letzten Kante. Hooks halten keine Lader-Sperre; ein Thread, der in einem Hook auf Nutzercode wartet, wartet nicht auf den Lader. Ein Kreis, an dem solche Kanten beteiligt sind, geht vom Nutzercode aus und ist nicht Sache des Laders. Ein vollständiger Beweis müsste zusätzlich Ausnahmen (ein ladender Thread, der abbricht, muss $z(k)$ zurücksetzen und die Wartenden wecken) und das Verhalten bei mehrfachen Abhängigkeiten auf denselben Schlüssel behandeln.
\end{proof}

\subsection{Auswirkungen auf die Sätze}

Die Charakterisierung des Erfolgs (Satz~\ref{satz:erfolg}) bleibt für jeden einzelnen Thread gültig; Lemma~\ref{lem:vertrag} wird zu einer Aussage über \emph{linearisierbare} Operationen, die Invariante von Satz~\ref{satz:refcount} wird unter $\mathcal{L}$ gehalten. Neu hinzu kommt eine Aussage über die Fairness des Wartens und eine Entscheidung, ob ein Thread beim Scheitern des ladenden Threads \emph{selbst} nochmals lädt oder den Fehler übernimmt; der Entwurf sollte den Fehler jedem Wartenden melden, damit das Ergebnis nicht von der Thread-Reihenfolge abhängt.

\section{Kanonischer Container und Signaturen}
\label{sec:signatur-ausblick}

Befund~B7 verlangt Echtheit. Der Entwurf von \texttt{proposal.md} nennt Hash und Signatur als Option (Abschnitt~10). Eine tragfähige Lösung hat drei Bausteine.

\subsection{Kanonische Kodierung}

\begin{definition}[Kanonischer Container, Entwurf]
	Ein Container heißt \emph{kanonisch}, wenn (1) hinter dem Payload-Bereich keine Bytes stehen, (2) die Tabelle nach Targets aufsteigend sortiert und ohne Doppelungen ist, (3) die Payload-Bereiche in Tabellenreihenfolge lückenlos und überlappungsfrei liegen, (4) die Metadaten als kanonisches JSON (feste Schlüsselreihenfolge, keine überflüssigen Leerzeichen, Zahlen in Normalform, etwa nach RFC~8785 \cite{rfc8785}) kodiert sind.
\end{definition}

\begin{satz}[Beweisskizze: Injektivität auf kanonischen Containern]
	Auf der Menge der kanonischen Container ist \texttt{unpack} injektiv, und $\mathrm{pack}\circ\mathrm{unpack}=\mathrm{id}$.
\end{satz}

\begin{proof}[Skizze]
	Aus (1) bis (3) folgt, dass die Längenangabe im Kopf, die Tabelle und die Offsets die Datei bis auf das letzte Byte bestimmen; (4) legt die Metadatenbytes fest. Die Kodierung durch \texttt{pack} erfüllt (1) bis (3) nach Lemma~\ref{lem:groesse}; sie erfüllt (4), wenn \texttt{json.dumps} durch die kanonische Kodierung ersetzt wird. Zwei kanonische Dateien mit gleichem Inhalt sind damit gleich.
\end{proof}

\subsection{Signierter Digest}

Gesucht ist ein Digest $D$, der Metadaten, Tabelle und Payload-Hashes bindet: $D=H(\mathrm{Kopf}\,\Vert\,\mathrm{Metadaten}\,\Vert\,\mathrm{Tabelle})$ mit SHA-256 \cite{fips180}; da die Tabelle die Hashes der Payloads enthält, bindet $D$ alle Payloads. Die Signatur $\mathrm{Sig}_{sk}(D)$ mit einem Verfahren wie Ed25519 \cite{rfc8032} wird in einen neuen Abschnitt hinter dem Payload-Bereich geschrieben, der von $D$ \emph{nicht} erfasst ist; der Kopf bekommt ein Feld für dessen Länge. Der Lader prüft die Signatur vor jeder Abbildung und gegen eine konfigurierte Menge vertrauter Schlüssel.

\begin{satz}[Beweisskizze: Sicherheit]
	Gibt es einen Angreifer, der eine akzeptierte Datei erzeugt, die nicht von einem Besitzer eines vertrauten Schlüssels stammt, so gibt es einen Angreifer gegen die Unfälschbarkeit des Signaturverfahrens oder gegen die Kollisionsresistenz von SHA-256.
\end{satz}

\begin{proof}[Skizze]
	Eine akzeptierte Datei hat eine gültige Signatur über $D'=H(\ldots)$. Ist $D'$ der Digest einer früher signierten Datei, so unterscheiden sich die Dateien in den gebundenen Teilen nicht (sonst Kollision) und in den Payloads nicht (die Hashes stehen in der Tabelle, wieder Kollision); sie sind bis auf ungebundene Teile gleich, die der Lader nicht auswertet. Andernfalls hat der Angreifer eine Signatur für einen neuen Wert $D'$ gefälscht.
\end{proof}

\subsection{Weitere Fragen}

Wie wird mit abgelaufenen oder zurückgezogenen Schlüsseln umgegangen (Sperrliste, Zeitstempel)? Wie wird die Signatur eines Fat-Containers behandelt, dem später ein Payload für ein neues Target hinzugefügt wird (Signatur über jeden Payload einzeln, mit Wurzel-Signatur)? Wie interagiert die Signatur mit reproduzierbaren Builds (gleiche Eingaben liefern gleiche Bytes)? Das sind Entwurfsfragen, die vor jeder Implementierung zu klären sind.

\section{Eine Grammatik für \texttt{SL\_PATH}}
\label{sec:slpath-ausblick}

Befund~B1 und die offene Frage des Entwurfs (Abschnitt~14) verlangen eine einzige Regel. Ein Vorschlag ist, \emph{nur} das Semikolon als Trennzeichen zu verwenden und \texttt{\textbackslash;} und \texttt{\textbackslash\textbackslash} als Maskierung. Der Doppelpunkt hat dann keine Sonderrolle, Laufwerksbuchstaben benötigen keinen Sonderfall, und Pfade mit Semikolon sind darstellbar.

\begin{satz}[Beweisskizze: Unbedingte Rundreise]
	Sei $\mathrm{esc}'$ die Maskierung von \texttt{\textbackslash} und \texttt{;} und $\mathrm{join}'$ die Verbindung mit \texttt{;}. Für jede Liste nichtleerer Zeichenketten gilt $\mathrm{split}'(\mathrm{join}'(p_1,\ldots,p_k))=(p_1,\ldots,p_k)$ ohne Randbedingung.
\end{satz}

\begin{proof}[Skizze]
	Wie im Beweis von Satz~\ref{satz:split}(c), mit \texttt{;} statt \texttt{:} als Trennzeichen. Der Laufwerkszweig entfällt, und mit ihm die Randbedingung. Eingaben mit einem einzelnen \texttt{\textbackslash} am Ende bleiben eine eigene Entscheidung: entweder Fehler oder wörtliches Zeichen; der Entwurf sollte sie festlegen und testen.
\end{proof}

Die Entscheidung hat einen Preis: Unter POSIX ist der Doppelpunkt der gewohnte Trenner; die Umstellung bricht Skripte, die \texttt{SL\_PATH} wie \texttt{PATH} setzen. Als Übergang bietet sich eine Warnung an, wenn die Zeichenkette einen Doppelpunkt enthält und keinen Laufwerksbuchstaben. Alternativ erlaubt eine Listen-API (der Entwurf nennt sie) die Umgehung der Zeichenkette ganz.

\section{Saubere Trennung zwischen Lader und Plattform}
\label{sec:adapter-ausblick}

Tabelle~\ref{tab:leaks} zeigt Plattformzugriffe außerhalb des Adapters. Sie lassen sich beseitigen, indem der Lader nicht mit Zeichenketten, sondern mit einer \emph{Plattform} arbeitet, die von außen übergeben wird:
\begin{itemize}
	\item einem Pfadtyp mit Operationen \texttt{join}, \texttt{parent}, \texttt{is\_absolute}, der für POSIX und Windows rein (ohne Dateisystemzugriff) und plattformunabhängig implementiert ist (in Python: \texttt{PurePosixPath}, \texttt{PureWindowsPath});
	\item einer Dateisystem-Schnittstelle \texttt{is\_file}, \texttt{realpath}, \texttt{mode\_of}, \texttt{owner\_of}, die der Adapter implementiert und die in Tests durch ein Fake-Dateisystem ersetzt wird;
	\item den Vorgaben für Plattform und Architektur als Teil des Adapters.
\end{itemize}
Mit dieser Trennung wird Satz~\ref{satz:parametrisch} zur Aussage ohne Ausnahmen: Die Beobachtung ist eine Funktion von $(E,R,A,h)$ und dem Orakel einschließlich Dateisystem. Als Gewinn fiele der Test \texttt{test\_hardened\_rejects\_world\_writable} aus dem plattformabhängigen Teil heraus: Er liefe mit einem Fake-Dateisystem auf jedem System gleich, und der Testname-Sonderfall in \texttt{\_world\_writable} (B3) entfiele.

\subsection{Ein Lader-Kern in C}

Der Entwurf nennt als Stufe~1 einen kleinen C-Kern mit Plattformadaptern. Er hat zwei Folgen. Erstens erhielten Sprachbindungen (Rust, Java, Node) \emph{denselben} Lader und nicht eine Neuimplementierung: Der Satz \ref{satz:parametrisch} gälte für den Kern, und jede Bindung wäre eine dünne Hülle. Zweitens entfiele die Abhängigkeit des Laders von der Python-Laufzeit; das Laden aus C++-Programmen (etwa \repo{csubgraph}) wäre direkt möglich. Der Preis ist ein zweiter Code, der gegen dieselbe Spezifikation geprüft werden muss; hierfür taugen sprachneutrale Prüfvektoren (Abschnitt~\ref{sec:konformitaet-ausblick}).

\subsection{Ein eigener Abbilder (Stufe~2)}

Die Abbildung über eine temporäre Datei (Listing~\ref{lst:adapter}) hat drei Schwächen: Sie braucht ein Verzeichnis mit Ausführungsrecht; sie hinterlässt kurzzeitig eine Datei, deren Schutz vom Betriebssystem abhängt; und sie überlässt dem Systemlader die Auflösung der Symbole zwischen Bibliotheken. Alternativen:
\begin{itemize}
	\item Anonyme Dateien im Speicher (unter Linux \texttt{memfd\_create} mit \texttt{dlopen} über den Dateideskriptor); nicht portabel, aber in einem Linux-Adapter möglich.
	\item Ein Abbilder im Benutzerraum, der ELF, PE und Mach-O selbst in den Speicher bringt, Relokationen anwendet und Symbole selbst auflöst. Das macht das Verhalten unabhängig von den Eigenheiten der Systemlader und löst B6, ist aber ein großer Eingriff und vergrößert die Angriffsfläche, wie der Entwurf in Abschnitt~8.5 selbst einräumt.
	\item Betriebssystemunterstützung (Stufe~3 des Entwurfs): politisch und technisch am aufwendigsten.
\end{itemize}

\section{Native Abhängigkeiten, Laufzeit und ABI-Entwicklung}
\label{sec:runtime-ausblick}

\subsection{Das Feld \texttt{runtime} mit Bedeutung füllen}

Das Feld hat die Werte \texttt{embedded} und \texttt{external}, wird aber nicht ausgewertet (B11). Eine mögliche Bedeutung: Bei \texttt{embedded} bringt der Payload seine C++-Laufzeit statisch mit (wie die statisch gelinkte MinGW-Laufzeit in Gen-DB); der Build prüft das (etwa über die Linker-Optionen), und der Lader muss nichts tun. Bei \texttt{external} wird die Laufzeit als eigene \SL-Abhängigkeit deklariert (\texttt{libstdcxx} mit Version), und der Lader löst sie wie jede Abhängigkeit auf. Beide Fälle ersetzen die Plattformsonderfälle aus Abschnitt~\ref{sec:evidenz}.

\subsection{Symbolauflösung zwischen Bibliotheken}

Für B6 sind drei Wege denkbar: der Modus \texttt{RTLD\_GLOBAL} für Abhängigkeiten, die Symbole bereitstellen, bei Beibehaltung des lokalen Modus für Blätter; eine Angabe je Abhängigkeit in den Metadaten (\texttt{"binding": "global"}); oder der eigene Abbilder. Der Test dazu ist einfach: Eine Basis exportiert eine Funktion, eine zweite Bibliothek ruft sie auf, ohne sie selbst zu enthalten, und \texttt{sl.open} der zweiten muss gelingen.

\subsection{Regeln der ABI-Entwicklung}

Der Lader prüft bisher $\alpha\in[\mathrm{low},\mathrm{high}]$ mit der Vorgabe $(1,1)$, also Gleichheit. Für eine Praxis mit mehreren Versionen braucht es Regeln: Wann erhöht sich $\alpha$? Gilt eine Bibliothek mit $\alpha=3$ als verträglich mit Nutzern, die $\alpha=2$ erwarten? Ein Vorschlag: $\alpha$ wird nur bei inkompatiblen Änderungen der exportierten C-Schnittstelle erhöht; Erweiterungen (neue Exporte) erhöhen nur die Minor-Version; jede Abhängigkeit trägt neben dem Versionsbereich einen ABI-Bereich. Die Bereichssprache (Satz~\ref{satz:konvex}) trägt das ohne Änderung, weil ABI-Bereiche Intervalle sind.

\section{Signaturen, Besitz und weitere Sprachen}

Die Signaturen der Metadaten sind C-Deklarationen und erlauben nur Skalare und Zeiger (Satz~\ref{satz:typen}). Für wirkliche Bibliotheken fehlen mindestens: Strukturen (\texttt{ctypes.Structure}) mit Layout-Prüfung, Rückrufe (Funktionszeiger), Felder fester Länge, Aufzählungen, Eigentums- und Lebensdaueranmerkungen (\emph{besitzt der Aufrufer den Rückgabezeiger, und mit welcher Funktion gibt er ihn frei?}), Fehlerrückgabe über Ausgabeparameter. Eine kleine Schnittstellenbeschreibung (IDL) in den Metadaten ersetzt dann die C-Deklaration als Quelle der Wahrheit; aus ihr erzeugt \texttt{sl gen} Header \emph{und} Bindings für Python, Rust (\texttt{libloading}), Java (Panama) und Node-API. Die Aussage von Satz~\ref{satz:schichtkonsistenz} (alle Schichten sind Funktionen der Metadaten) bleibt der Leitgedanke; sie muss für jede neue Sprache bewiesen oder getestet werden. Zu klären ist, wie das Datenmodell (LLP64/LP64) in der IDL ausgedrückt wird: Die Empfehlung ist, nur Typen fester Breite zuzulassen; dann entfiele der Datenmodellunterschied an der Schnittstelle ganz.

\section{Build-System und Werkzeugkette}

Nach B9 und B10 gehören zum Ausbau: Compiler- und Standardwahl, Linker-Angaben, Beschreibung von GNU- und MSVC-Toolchains als Daten (nicht als \texttt{if}-Zweige), ein belastbares CMake-Modul, Querübersetzung für die Zielmatrix, und ein Befehl \texttt{sl doctor}, der Compiler, Zielplattformen, Schreibrechte der Suchverzeichnisse und das Temporärverzeichnis prüft und das Ergebnis lesbar ausgibt. Die Build-Matrix (Satz~\ref{satz:schichtkonsistenz}) wird dann nicht nur berechnet, sondern in CI auf allen drei Systemen ausgeführt, mit je einer Zeile je Target, einer Zusage je Zeile und den Kosten pro Lauf. Reproduzierbare Builds sind Voraussetzung für Signaturen (Abschnitt~\ref{sec:signatur-ausblick}).

\section{Versionen: Diagnose statt Suche}

Satz~\ref{satz:npvoll} schließt aus, dass der Lader unter Kandidaten sucht; er schließt \emph{nicht} aus, dass er bei Mehrdeutigkeit hilft. Zwei Wege erhalten die Polynomialzeit:

\begin{enumerate}
	\item \emph{Diagnose.} Beim Scheitern an einem Bereich listet der Lader alle weiter hinten gefundenen Kandidaten mit Versionen auf (ein Durchlauf über die Kandidatenliste, $O(n_{\max})$ Dateikopf-Lesungen). Die Fehlermeldung nennt dann, dass \texttt{d2/base.sl} (Version 2.0.0) den Bereich erfüllt, aber hinter \texttt{d1/base.sl} liegt (Satz~\ref{satz:unvollstaendig}(a)).
	\item \emph{Sperrdatei.} Die Auswahl wird außerhalb des Laders getroffen (Paketverwaltung oder Build), und das Ergebnis, eine Abbildung Name $\to$ (Pfad, Hash), liegt als Sperrdatei vor. Der Lader \emph{prüft} sie nur. Das passt zu Satz~\ref{satz:npvoll}: Eine zulässige Auswahl ist ein Zeugnis in NP, und seine Prüfung ist polynomiell. Die Sperrdatei ist also genau das Zertifikat, das ein Lader in Polynomialzeit prüfen kann, auch wenn das Finden hart ist \cite{cox2016,abate2012}.
\end{enumerate}

Zu den Metadaten gehört außerdem die Prüfung, dass Bereiche erfüllbar sind (Satz~\ref{satz:konvex}(d)); sie ist billig, weil die Erfüllbarkeit einer Bereichsliste durch Vergleich der größten unteren mit der kleinsten oberen Schranke entschieden wird.

\section{Konformitätssuite mit sprachneutralen Vektoren}
\label{sec:konformitaet-ausblick}

Der Entwurf verlangt eine Konformitätssuite (Abschnitt~5.9), die jeder Lader, Adapter und jede Bindung bestehen muss. Die Tests der Referenzimplementierung sind dafür nicht geeignet, weil sie Python-Aufrufe prüfen. Sprachneutral sind \emph{Prüfvektoren}: Eine Datei beschreibt ein Dateisystem (Pfade, Inhalte, Rechte), eine Umgebung, eine Folge von Operationen und die erwartete Beobachtung (Statuscodes, \texttt{failed\_name}, \texttt{searched}, Reihenfolge der Hook-Aufrufe). Die Beobachtung ist nach Definition~\ref{def:parametrizitaet} genau das, worauf der Entwurf Wert legt. Aus dem Python-Modell können die Vektoren erzeugt werden, einschließlich der Zufallsfälle aus Abschnitt~\ref{sec:verifikation-ausblick}; ein Lader in C oder Rust wird gegen sie geprüft, ohne dass der Python-Code dafür vorhanden sein muss. Die Vektoren zu Satz~\ref{satz:split} (Zerlegung von \texttt{SL\_PATH}) und Satz~\ref{satz:struktur} (Kandidatenliste) sind besonders einfach, weil sie reine Funktionen prüfen.

\section{Migration der Beispielprojekte}
\label{sec:migration}

Die Wirkung von \SL\ auf Gen-DB und \texttt{subgraph-sat-solver-econ} ist bislang eine Behauptung des Entwurfs. Ein Plan zum Prüfen:

\begin{enumerate}
	\item \emph{Metadaten schreiben.} Für die C-Schnittstelle von \texttt{csubgraph} (Shim für Gen-DB) und für \texttt{wsat\_shim} je eine Metadatendatei mit Name, Version, ABI-Version, Targets, Exporten, Signaturen und Hooks. Die Signaturen sind vorhanden (als C-Header); sie müssen in das Fragment aus Definition~\ref{def:fragment} passen. Der Quelltext dieser Projekte lag der Arbeit nicht vor.
	\item \emph{Bau mit \texttt{sl build}.} Voraussetzungen sind die Lücken aus B9 (C++17, Linker, \texttt{cl}).
	\item \emph{Laden mit \texttt{sl.open}} anstelle von \texttt{ctypes.CDLL}; die alte Pfadlogik bleibt zunächst als Rückfall.
	\item \emph{Messung}: Wie viele Zeilen Plattformcode entfallen (Zählung vor und nach), wie viele Fehlerfälle werden eindeutig gemeldet, welche Zeit kostet das Laden, wie ändert sich die README (Abschnitte, die entfallen).
	\item \emph{Entscheidung nach Messung}, wie der Entwurf selbst es in seinem Fahrplan vorsieht (Abschnitt~13).
\end{enumerate}

Als Hypothesen sind vorab festzulegen: \textbf{H1:} Die Zahl der Plattformverzweigungen in den Ladepfaden beider Projekte sinkt auf null. \textbf{H2:} Die Ladezeit steigt um weniger als eine feste Schranke (etwa $50\,\mathrm{ms}$ je Bibliothek für Payloads unter $10\,\mathrm{MB}$). \textbf{H3:} Jeder in Abschnitt~\ref{sec:evidenz} genannte Fehlerfall liefert einen Statuscode des Fehlermodells statt einer Systemmeldung. Jede Hypothese ist falsifizierbar; \offen{offen} bis zur Messung.

\section{Kostenmodell und Messungen}

Die Kosten von \texttt{open} setzen sich zusammen aus: Suchen ($\le n_{\max}$ Dateisystemprüfungen), Lesen der Datei (linear in der Dateigröße $S$), Entpacken (der Code kopiert den Rest der Datei in \texttt{blob}, weitere $O(S)$), Hashen des Payloads (linear), Schreiben der temporären Datei (linear) und \texttt{dlopen}. Als Modell
\[
T_{\mathrm{open}}(S)\approx a+b\,S,\qquad b=\tfrac{1}{v_{\mathrm{Lesen}}}+\tfrac{1}{v_{\mathrm{Kopie}}}+\tfrac{1}{v_{\mathrm{SHA}}}+\tfrac{1}{v_{\mathrm{Schreiben}}},
\]
mit einer Konstante $a$ für Suche und \texttt{dlopen}. Zu messen sind $a$ und $b$ je System, mit kalten und warmen Zwischenspeichern, und der Vergleich mit direktem \texttt{ctypes.CDLL}. Bei den hier betrachteten Anwendungen (wenige Aufrufe mit großen Datenmengen, Entwurf Abschnitt~11) wird der Lader-Aufwand voraussichtlich klein gegen die Rechenzeit sein; ob das stimmt, ist \offen{offen}. Verbesserungen, die das Modell nahelegt: Lesen mit Speicherabbildung statt Kopie, Hash nur beim ersten Laden und zwischengespeichert nach Dateistempel, Wiederverwendung der materialisierten Datei (Inhaltsadresse) statt jedes Mal neu zu schreiben, was zugleich die Temporärverzeichnisse spart.

\section{Verbindung zum Subgraph-Programm}

Für die Reihe der Arbeiten zum Subgraph Algorithmus ergeben sich drei Anschlüsse.
\begin{enumerate}
	\item \emph{Katalog.} Der Katalog der Wiederverwendungen (\W{1} bis \W{18}) in \texttt{csubgraph.tex} sollte um einen Eintrag der Form \emph{Infrastruktur} ergänzt werden: das \SL-Format als Hülle der Bibliothek (Abschnitt~\ref{sec:formen}). Die Zuordnung zu den Formen wäre unverändert.
	\item \emph{Orakel als Bibliothek.} Der Subgraph-SAT-Solver ruft die Subgraph-Bibliothek als Orakel (Vorarbeit). Ein \SL-Container mit einem Export, der ein Graphenpaar entgegennimmt und eine Einbettung zurückgibt, wäre eine Schnittstelle, hinter der die Python-Referenz, die C++-Bibliothek und künftige Implementierungen austauschbar sind. Satz~\ref{satz:schichtkonsistenz} beschreibt, wie die Signatur dafür zur einzigen Quelle würde.
	\item \emph{Beweisstil.} Die Kette 3-SAT $\to$ Clique $\to$ Teilgraph und die Reduktion in Satz~\ref{satz:npvoll} sind verwandt; die Reduktion auf die Versionsauswahl könnte in der Reihe als Beispiel dafür dienen, dass die Schwere eines Problems die Gestalt der Schnittstelle begründet (hier: der Lader wählt nicht).
\end{enumerate}

\section{Standardisierung und Betrieb}

Der Entwurf stellt die Frage, wer Spezifikation, Referenzimplementierung und Konformitätssuite pflegt (Abschnitt~14). Ein tragfähiger Weg ist, die Spezifikation als versioniertes Dokument mit nummerierten Anforderungen zu führen, wie es Tabelle~\ref{tab:konformitaet} vorbereitet, jede Anforderung an mindestens einen Prüfvektor zu binden und die Abweichungen der Referenzimplementierung (B14) offen zu führen. Das Risiko der Zersplitterung (Entwurf Abschnitt~11, \emph{xkcd~927}) wird durch den Legacy-Modus gemildert, der gewöhnliche \texttt{.dll}-, \texttt{.so}- und \texttt{.dylib}-Dateien über dieselbe API öffnet; er fehlt bisher (B14(g)) und ist für die Übernahme in bestehende Projekte vermutlich wichtiger als jede andere Erweiterung. Zur Frage der Lizenz und des Betriebs: Das Repository enthält eine \texttt{LICENSE}-Datei; für eine Normierung wäre eine Lizenz zu wählen, die Übernahme in Systembibliotheken nicht verbietet.

\section{Weitere Entwürfe des Entwurfs}

Der Entwurf nennt Fragen, die diese Arbeit nicht beantwortet, aber präzisiert:
\begin{itemize}
	\item \emph{Mehrere Architekturen in einer Datei} (x86-64, ARM64): bereits durch den Fat-Container gelöst (Target = Plattform + Architektur); offen ist die Auswahl zwischen verträglichen Targets (etwa \texttt{x86} auf \texttt{x86\_64}) und die Rolle von \texttt{musl} gegenüber \texttt{glibc}, die das Target nicht ausdrückt.
	\item \emph{Verzögertes Binden und Laden:} Welche Semantik gilt für Symbole, die erst bei Gebrauch aufgelöst werden? Im Modell wäre ein dritter Zustand je Symbol nötig; Satz~\ref{satz:export} bliebe wahr, aber Fehler träten später auf.
	\item \emph{Entladen von Bibliotheken mit Threads oder Rückrufen:} Bemerkung zu \texttt{dlclose} in Kapitel~\ref{chap:adapter}; eine Zusage braucht einen Mechanismus, der Threads der Bibliothek beendet (ein \texttt{fini}-Vertrag).
	\item \emph{Einbettung des Laders in Prozesse ohne Ausführungsrecht im Temporärverzeichnis:} Hinweis zur Stufe~2.
\end{itemize}

\section{Offene Forschungsfragen}

\begin{enumerate}
	\item Lässt sich die Menge $\{$erfolgreiche Abläufe$\}$ des Laders mit Sperren je Schlüssel so charakterisieren, dass Satz~\ref{satz:erfolg} unverändert für jeden Thread gilt, und wie ist die Fairness des Wartens zu fassen?
	\item Gibt es eine \emph{minimale} Spezifikation der Ladeprozedur, aus der sich alle Sätze dieser Arbeit ableiten lassen, und welche Annahmen sind darin unvermeidbar?
	\item Welche Sicherheitseigenschaft gilt für den gehärteten Modus gegen den Angreifer (A3) formal, wenn Besitzer und Rechte aller übergeordneten Verzeichnisse geprüft werden, und was bleibt als nicht entscheidbarer Rest (etwa Zugriffsrechte, die sich zwischen Prüfung und Lesen ändern)?
	\item Ist \textsc{DepSelect} auch NP-vollständig, wenn jede Version höchstens eine Anforderung hat und jedes Paket nur zwei Versionen? (Die Reduktion in Satz~\ref{satz:npvoll} braucht drei Versionen für Klauselpakete; für zwei ist die Frage offen und hängt mit 2-SAT zusammen.)
	\item Welche Klasse von Signaturen lässt sich durch ein \emph{beweisbar} ABI-treues Binding erfassen, und wie viel von der C-ABI der drei Plattformen muss das Modell kennen?
	\item Wie groß ist der Anteil des Plattformcodes in Projekten, die C/C++-Bibliotheken aus Python nutzen, tatsächlich? Der Entwurf behauptet ``tausendfach wiederholt''; eine Stichprobe öffentlicher Projekte würde die Größenordnung belegen oder widerlegen.
	\item Wie verhält sich das Format zu WebAssembly mit Komponentenmodell, das ein Format und mehrere Laufzeiten verspricht: Welche Eigenschaften des Entwurfs (Export, Fehlermodell, Lebenszyklus) lassen sich dort ausdrücken, und was kostet die Schnittstelle?
\end{enumerate}

